\chapter{Justificación}
La justificación de este proyecto surge de la necesidad de fortalecer la investigación en robótica dentro de la Universidad Santo Tomás y de aprovechar plenamente la plataforma cuadrúpeda disponible en el laboratorio de robótica, la cual actualmente no se encuentra habilitada para su utilización en docencia, investigación y extensión universitaria debido a su estado de construcción y a diversas limitaciones técnicas. Reactivar esta plataforma representa una oportunidad estratégica para impulsar la generación de conocimiento en áreas de locomoción robótica, inteligencia artificial y control de sistemas dinámicos, campos fundamentales para la formación de ingenieros en un entorno tecnológico en constante evolución. En este contexto, en Colombia existe un creciente impulso hacia la automatización y el desarrollo de tecnologías inteligentes, lo que refuerza la pertinencia de este tipo de iniciativas académicas y de investigación.\

\textbf{Científica:} Este trabajo contribuye directamente al avance del conocimiento en locomoción cuadrúpeda, permitiendo que docentes y estudiantes dispongan de un recurso experimental funcional para desarrollar proyectos, validar teorías y fortalecer competencias prácticas en robótica.\

\textbf{Tecnológica:} Convertir esta plataforma en un sistema operativo y completamente caracterizado permitirá implementar control, algoritmos y sistemas robóticos funcionales, abriendo la puerta a nuevos proyectos de investigación, acciones de mejora para trabajos de grado y desarrollos tecnológicos que incrementan el valor del laboratorio y enriquecen la producción académica institucional.\

\textbf{Académica:} La reactivación de la plataforma impactará de manera directa varios espacios académicos del programa, particularmente aquellos relacionados con control automático, instrumentación electrónica, procesamiento y adquisición de datos, robótica, sistemas embebidos y asignaturas del nuevo plan curricular orientadas a la automatización inteligente. Además, la iniciativa se articula con los Objetivos de Desarrollo Sostenible (ODS 4, 8 y 9), fortaleciendo procesos formativos mediante tecnologías avanzadas, impulsando la investigación e innovación en sistemas robóticos y fomentando competencias tecnológicas que responden a las demandas de los sectores productivos.\

En conjunto, este proyecto se justifica porque promueve la investigación universitaria, fortalece la formación en robótica e inteligencia artificial, impulsa el avance tecnológico regional, articula espacios académicos clave del programa y contribuye al desarrollo sostenible. Reactivar la plataforma cuadrúpeda no solo resuelve una necesidad inmediata del laboratorio, sino que habilita nuevas oportunidades académicas, investigativas y formativas de alto impacto.


